\tzpanchor{0153}
\tzpentryheader{0153}{Gravitational Stress--Energy Contribution}

\begingroup\small
\begingroup\scriptsize
\setlength{\tabcolsep}{4pt}
\renewcommand{\arraystretch}{0.92}
\noindent\begin{tabularx}{\linewidth}{@{}p{0.24\linewidth}p{0.36\linewidth}X@{}}
\textbf{UUID} & \multicolumn{2}{l@{}}{\tzpcode{6fb08435-6667-5e5c-b92c-e572d0bdfc47}}\\
\textbf{PiID} & \tzpcode{1174502841027} & \textbf{TZPi}\quad \tzpcode{2}\quad \textbf{PiPE}\quad \tzpcode{160}\\
\textbf{RTEID} & \textbf{Poloidal $\theta$}\quad \tzpcode{594.1330 rad} & \textbf{Toroidal $\phi$}\quad \tzpcode{00.9789 rad}\\
\textbf{TZPID} & \tzpcode{ID0153} & \textbf{Node Dependencies}\quad \tzpidref{0152}\\
\textbf{Registry} & \tzpcode{registry\_id\_0153} & \textbf{Scale}\quad curvature-sector closure term\\
\textbf{LOC Classification} & QC178 Einstein equations & QA767 differential geometry\\
\textbf{arXiv Classification} & Primary: gr-qc \quad Cross-list: hep-th, math-ph & \\
\textbf{Primary Support} & Curvature Package & Einstein Tensor\\
\textbf{Closure Support} & Einstein-Form Rewriting & \\
\end{tabularx}\par
\endgroup
\endgroup

\tzpsection{Canonical Statement}
The gravitational sector is written as the curvature-side source term that matches the stress bookkeeping used for the vacuum and matter sectors. It packages Ricci curvature, scalar curvature, and the cosmological term into one tensor contribution.

\tzpsection{Canonical Equation}
\[
T_{\mathrm{grav}}^{\mu\nu}
=
\frac{c^4}{8\pi G}
\left(
R^{\mu\nu}-\frac{1}{2}g^{\mu\nu}R+\Lambda g^{\mu\nu}
\right)
\]
\[
G_{\mu\nu}=R_{\mu\nu}-\frac{1}{2}g_{\mu\nu}R
\]
\[
T_{\mathrm{grav}}^{\mu\nu}
=
\frac{c^4}{8\pi G}
\left(
G^{\mu\nu}+\Lambda g^{\mu\nu}
\right)
\]

\begin{minipage}[t]{0.49\linewidth}
\tzpsection{Canonical Symbol Key}
\begingroup
\small
\raggedright
\textbf{Source equation symbols.}\par
\begin{itemize}[leftmargin=1.35em,itemsep=1pt,topsep=1pt,parsep=0pt]
    \item $T_{\mathrm{grav}}$: source equation symbol.
    \item $\mu$: source equation symbol.
    \item $c$: speed of light or relativistic conversion factor.
    \item $\pi$: source equation symbol.
    \item $G$: source equation symbol.
    \item $R$: global radius, boundary radius, or topological scale.
    \item $\Lambda$: lemniscatic representative or canonical branch object.
    \item $G_{\mu\nu}$: source equation symbol.
    \item $R_{\mu\nu}$: source equation symbol.
\end{itemize}
\endgroup
\end{minipage}\hfill
\begin{minipage}[t]{0.47\linewidth}
\tzpsection{Wolfram Form}
\begingroup
\scriptsize
\raggedright
\textbf{Source Wolfram model.}\par
\begin{Verbatim}[fontsize=\tiny,breaklines=true,breakanywhere=true]
(* TZPID ID0153 generated source Wolfram model *)
ClearAll[K0153, Cr0153, T, R, A, W, I, N];
eq0153$1 = HoldForm[TGrav^munu == (c^4/(8*Pi*G))*(R^munu - (1/2)*g^munu*R + Lambda*g^munu)];
eq0153$2 = HoldForm[GMunu == RMunu - (1/2)*gMunuR];
eq0153$3 = HoldForm[TGrav^munu == (c^4/(8*Pi*G))*(G^munu + Lambda*g^munu)];

K0153 = HoldForm[{eq0153$1, eq0153$2, eq0153$3}];

N = "normalized TRAWIN closure object for Gravitational Stress--Energy Contribution";
I = "Einstein-Form Rewriting integration/comparison layer";
W = "Einstein Tensor wave or operator carrier";
A = "Curvature Package admissible action/amplitude condition";
R = "Einstein Tensor rotation/curl preservation layer";
T = "Curvature Package local source condition";

Cr0153[expr_] := HoldForm[N[I[W[A[R[T[expr]]]]]]];

Result0153 = Cr0153[K0153];
\end{Verbatim}
\endgroup
\end{minipage}

\par\vspace{0.45\baselineskip}
\noindent\begin{minipage}[t]{\linewidth}
\begingroup
\small
\raggedright
\textbf{TRAWIN Operator Key.}\par
\noindent\textbf{Time/localization operator:} $\mathsf{T}\!\left[\mathrm{local}\right]$ Curvature Package local source condition.\par
\noindent\textbf{Rotation/curl operator:} $\mathsf{R}\!\left[\nabla\times\right]$ Einstein Tensor rotation/curl preservation layer.\par
\noindent\textbf{Action/amplitude operator:} $\mathsf{A}\!\left[\mathrm{admissible}\right]$ Curvature Package admissible action/amplitude condition.\par
\noindent\textbf{Wave/d'Alembertian operator:} $\mathsf{W}\!\left[\Box\right]$ Einstein Tensor wave or operator carrier.\par
\noindent\textbf{Integration operator:} $\mathsf{I}\!\left[\int\right]$ Einstein-Form Rewriting integration/comparison layer.\par
\noindent\textbf{Normalization operator:} $\mathsf{N}\!\left[\mathrm{normalized}\right]$ normalized TRAWIN closure object for Gravitational Stress--Energy Contribution.\par
\endgroup
\end{minipage}

\clearpage
\noindent\textbf{[ID 0153] Gravitational Stress--Energy Contribution}\par
\vspace{0.35em}
\tzpsection{Derivation Layer}
\paragraph{Primary Equation.}
\[
T_{\mathrm{grav}}^{\mu\nu}
=
\frac{c^4}{8\pi G}
\left(
R^{\mu\nu}-\frac{1}{2}g^{\mu\nu}R+\Lambda g^{\mu\nu}
\right)
\]
\[
G_{\mu\nu}=R_{\mu\nu}-\frac{1}{2}g_{\mu\nu}R
\]
\[
T_{\mathrm{grav}}^{\mu\nu}
=
\frac{c^4}{8\pi G}
\left(
G^{\mu\nu}+\Lambda g^{\mu\nu}
\right)
\]

\paragraph{Primary Derivation.}
The relation is obtained by differentiating and applying the relevant definitions. leading to the displayed result.
\paragraph{Equation Type.}
This statement is classified as a other relation.

\paragraph{Interpretation.}
This relation applies to the quantities defined above. Interpreted through TRAWIN, the equation functions as the generalized carrier for the entry’s information content. According to CHL, the derivation provides a witness morphism connecting the canonical proposition to its proof certificate.

\par\vspace{0.35em}
\noindent\textbf{Derivable Consequences.}\quad
\begingroup
\footnotesize
\raggedright
The canonical equation anchors Gravitational Stress--Energy Contribution by connecting the source condition to the derivation layer and proof certificate.
\par\endgroup

\tzpsection{Equation Support}
\noindent\textbf{1. Curvature Package}\par
\[
T_{\mathrm{grav}}^{\mu\nu}
=
\frac{c^4}{8\pi G}
\left(
R^{\mu\nu}-\frac{1}{2}g^{\mu\nu}R+\Lambda g^{\mu\nu}
\right)
\]
\noindent This defines the gravitational sector directly from curvature data.\par
\noindent\textbf{2. Einstein Tensor}\par
\[
G_{\mu\nu}=R_{\mu\nu}-\frac{1}{2}g_{\mu\nu}R
\]
\noindent This isolates the standard curvature combination used in the closure law.\par
\noindent\textbf{3. Einstein-Form Rewriting}\par
\[
T_{\mathrm{grav}}^{\mu\nu}
=
\frac{c^4}{8\pi G}
\left(
G^{\mu\nu}+\Lambda g^{\mu\nu}
\right)
\]
\noindent This closes the progression by rewriting the gravitational contribution in Einstein-tensor form.\par

\clearpage
\noindent\textbf{[ID 0153] Gravitational Stress--Energy Contribution}\par
\vspace{0.35em}
\tzpsection{Dictionary}
\begingroup\small
\tzpsection{Definition}
Gravitational Stress--Energy Contribution is a canonical-registry. Matter / gyro contribution (generic imperfect fluid ansatz) For completeness include a general matter stress--energy contribution to be added to the EM piece with \$h\textasciicircum{}\{\textbackslash{}mu\textbackslash{}nu\}=g\textasciicircum{}\{\textbackslash{}mu\textbackslash{}nu\}+u\textasciicircum{}\{\textbackslash{}mu\}u\textasciicircum{}\{\textbackslash{}nu\}\$ and constraints \$u\_\textbackslash{}mu q\textasciicircum{}\textbackslash{}mu =0\$, \$u\_\textbackslash{}mu\textbackslash{}pi\textasciicircum{}\{\textbackslash{}mu\textbackslash{}nu\}=0\$, \$\textbackslash{}pi\textasciicircum{}\{\textbackslash{}mu\}\{\}\_\{\textbackslash{}mu\}=0\$. Use your gyro 4-velocity \$u\textasciicircum{}\textbackslash{}mu\$ from your notes to build this part. Eigen-decomposition and principal pressures To find principal pressures and principal axes, perform the eigen-decomposition of the spatial tensor \$S\textasciicircum{}i\{\}\_j\$: 1.

% BEGIN TZP CONTEXTUAL MEANING
\tzpsection{Contextual Meaning}
\begingroup\footnotesize\setlength{\emergencystretch}{3em}\sloppy
\textbf{Formal statement:} Matter / gyro contribution (generic imperfect fluid ansatz) For completeness include a general matter stress--energy contribution to be added to the EM piece with \$h\textasciicircum{}\{\textbackslash{}mu\textbackslash{}nu\}=g\textasciicircum{}\{\textbackslash{}mu\textbackslash{}nu\}+u\textasciicircum{}\{\textbackslash{}mu\}u\textasciicircum{}\{\textbackslash{}nu\}\$ and constraints \$u\_\textbackslash{}mu q\textasciicircum{}\textbackslash{}mu =0\$, \$u\_\textbackslash{}mu\textbackslash{}pi\textasciicircum{}\{\textbackslash{}mu\textbackslash{}nu\}=0\$, \$\textbackslash{}pi\textasciicircum{}\{\textbackslash{}mu\}\{\}\_\{\textbackslash{}mu\}=0\$. Use your gyro 4-velocity \$u\textasciicircum{}\textbackslash{}mu\$ from your notes to build this part. Eigen-decomposition and principal pressures To find principal pressures and principal axes, perform the eigen-decomposition of the spatial tensor \$S\textasciicircum{}i\{\}\_j\$: 1.
\textbf{Contextual meaning:} / Test Handle**: How this entry could be empirically challenged --> ID 150: Stress--Energy Conservation Constraint Type: Conservation Law Formal Statement: The total stress--energy tensor Ttotal  =Tvacuum  +Tmatter  +Tgravitational  T\_\{\textbackslash{}text\{total\}\}\textasciicircum{}\{\textbackslash{}mu\textbackslash{}nu\} = T\_\{\textbackslash{}text\{vacuum\}\}\textasciicircum{}\{\textbackslash{}mu\textbackslash{}nu\} + T\_\{\textbackslash{}text\{matter\}\}\textasciicircum{}\{\textbackslash{}mu\textbackslash{}nu\} + T\_\{\textbackslash{}text\{gravitational\}\}\textasciicircum{}\{\textbackslash{}mu\textbackslash{}nu\}Ttotal   =Tvacuum   +Tmatter   +Tgravitational   satisfies   Ttotal  =0\textbackslash{}nabla\_\{\textbackslash{}nu\} T\_\{\textbackslash{}text\{total\}\}\textasciicircum{}\{\textbackslash{}mu\textbackslash{}nu\} = 0   Ttotal   =0 at all points, with Tvacuum= TZP T   TZP exp (-r2/ coherence2)T\_\{\textbackslash{}text\{vacuum\}\} = \textbackslash{}langle \textbackslash{}text\{TZP\}|T\textasciicircum{}\{\textbackslash{}mu\textbackslash{}nu\}|\textbackslash{}text\{TZP\}\textbackslash{}rangle \textbackslash{}exp(-r 2/\textbackslash{}xi\_\{\textbackslash{}text\{coherence\}\} 2)Tvacuum = TZP T   TZP exp(-r2/ coherence2 ), TmatterT\_\{\textbackslash{}text\{matter\}\}Tmatter defined from effective energy density and pressure, ...
\textbf{Derivable consequences:} Source evidence records the entry as a reusable object whose consequences should be read from the matched formal statement and local proof route.
\textbf{Lemma support:} Matter / gyro contribution (generic imperfect fluid ansatz) For completeness include a general matter stress--energy contribution to be added to the EM piece with \$h\textasciicircum{}\{\textbackslash{}mu\textbackslash{}nu\}=g\textasciicircum{}\{\textbackslash{}mu\textbackslash{}nu\}+u\textasciicircum{}\{\textbackslash{}mu\}u\textasciicircum{}\{\textbackslash{}nu\}\$ and constraints \$u\_\textbackslash{}mu q\textasciicircum{}\textbackslash{}mu =0\$, \$u\_\textbackslash{}mu\textbackslash{}pi\textasciicircum{}\{\textbackslash{}mu\textbackslash{}nu\}=0\$, \$\textbackslash{}pi\textasciicircum{}\{\textbackslash{}mu\}\{\}\_\{\textbackslash{}mu\}=0\$. Use your gyro 4-velocity \$u\textasciicircum{}\textbackslash{}mu\$ from your notes to build this part. Eigen-decomposition and principal pressures To find principal pressures and principal axes, perform the eigen-decomposition of the spatial tensor \$S\textasciicircum{}i\{\}\_j\$: 1.\par
\endgroup
% END TZP CONTEXTUAL MEANING

\tzpsection{Technical Interpretation}
/ Test Handle**: How this entry could be empirically challenged --> ID 150: Stress--Energy Conservation Constraint Type: Conservation Law Formal Statement: The total stress--energy tensor Ttotal  =Tvacuum  +Tmatter  +Tgravitational  T\_\{\textbackslash{}text\{total\}\}\textasciicircum{}\{\textbackslash{}mu\textbackslash{}nu\} = T\_\{\textbackslash{}text\{vacuum\}\}\textasciicircum{}\{\textbackslash{}mu\textbackslash{}nu\} + T\_\{\textbackslash{}text\{matter\}\}\textasciicircum{}\{\textbackslash{}mu\textbackslash{}nu\} + T\_\{\textbackslash{}text\{gravitational\}\}\textasciicircum{}\{\textbackslash{}mu\textbackslash{}nu\}Ttotal   =Tvacuum   +Tmatter   +Tgravitational   satisfies   Ttotal  =0\textbackslash{}nabla\_\{\textbackslash{}nu\} T\_\{\textbackslash{}text\{total\}\}\textasciicircum{}\{\textbackslash{}mu\textbackslash{}nu\} = 0   Ttotal   =0 at all points, with Tvacuum= TZP T   TZP exp (-r2/ coherence2)T\_\{\textbackslash{}text\{vacuum\}\} = \textbackslash{}langle \textbackslash{}text\{TZP\}|T\textasciicircum{}\{\textbackslash{}mu\textbackslash{}nu\}|\textbackslash{}text\{TZP\}\textbackslash{}rangle \textbackslash{}exp(-r 2/\textbackslash{}xi\_\{\textbackslash{}text\{coherence\}\} 2)Tvacuum = TZP T   TZP exp(-r2/ coherence2 ), TmatterT\_\{\textbackslash{}text\{matter\}\}Tmatter defined from effective energy density and pressure, ...

\tzpsection{Equation Grounding}
Matter / gyro contribution (generic imperfect fluid ansatz) For completeness include a general matter stress--energy contribution to be added to the EM piece with \$h\textasciicircum{}\{\textbackslash{}mu\textbackslash{}nu\}=g\textasciicircum{}\{\textbackslash{}mu\textbackslash{}nu\}+u\textasciicircum{}\{\textbackslash{}mu\}u\textasciicircum{}\{\textbackslash{}nu\}\$ and constraints \$u\_\textbackslash{}mu q\textasciicircum{}\textbackslash{}mu =0\$, \$u\_\textbackslash{}mu\textbackslash{}pi\textasciicircum{}\{\textbackslash{}mu\textbackslash{}nu\}=0\$, \$\textbackslash{}pi\textasciicircum{}\{\textbackslash{}mu\}\{\}\_\{\textbackslash{}mu\}=0\$. Use your gyro 4-velocity \$u\textasciicircum{}\textbackslash{}mu\$ from your notes to build this part. Eigen-decomposition and principal pressures To find principal pressures and principal axes, perform the eigen-decomposition of the spatial tensor \$S\textasciicircum{}i\{\}\_j\$: 1.

\tzpsection{Interpretive Role}
/ Test Handle**: How this entry could be empirically challenged --> ID 150: Stress--Energy Conservation Constraint Type: Conservation Law Formal Statement: The total stress--energy tensor Ttotal  =Tvacuum  +Tmatter  +Tgravitational  T\_\{\textbackslash{}text\{total\}\}\textasciicircum{}\{\textbackslash{}mu\textbackslash{}nu\} = T\_\{\textbackslash{}text\{vacuum\}\}\textasciicircum{}\{\textbackslash{}mu\textbackslash{}nu\} + T\_\{\textbackslash{}text\{matter\}\}\textasciicircum{}\{\textbackslash{}mu\textbackslash{}nu\} + T\_\{\textbackslash{}text\{gravitational\}\}\textasciicircum{}\{\textbackslash{}mu\textbackslash{}nu\}Ttotal   =Tvacuum   +Tmatter   +Tgravitational   satisfies   Ttotal  =0\textbackslash{}nabla\_\{\textbackslash{}nu\} T\_\{\textbackslash{}text\{total\}\}\textasciicircum{}\{\textbackslash{}mu\textbackslash{}nu\} = 0   Ttotal   =0 at all points, with Tvacuum= TZP T   TZP exp (-r2/ coherence2)T\_\{\textbackslash{}text\{vacuum\}\} = \textbackslash{}langle \textbackslash{}text\{TZP\}|T\textasciicircum{}\{\textbackslash{}mu\textbackslash{}nu\}|\textbackslash{}text\{TZP\}\textbackslash{}rangle \textbackslash{}exp(-r 2/\textbackslash{}xi\_\{\textbackslash{}text\{coherence\}\} 2)Tvacuum = TZP T   TZP exp(-r2/ coherence2 ), TmatterT\_\{\textbackslash{}text\{matter\}\}Tmatter defined from effective energy density and pressure, ...
\endgroup

\tzpsection{TZP Thesaurus}
\begin{enumerate}[leftmargin=1.6em,itemsep=6pt,topsep=4pt]
\item \textbf{Curvature Package}: The definition of the gravitational sector directly from curvature data.
\item \textbf{Einstein Tensor}: The isolation of the standard curvature combination used in the closure law.
\item \textbf{Einstein-Form Rewriting}: The rewriting of the gravitational contribution in Einstein-tensor form.
\end{enumerate}

\tzpsection{Encyclopedia Mathematica}
\begingroup\footnotesize\sloppy
The visual plate below is reserved for the source-specific equation and progression graphic for [ID 0153]. It is included only as a companion to the canonical equation, not as a substitute for the formal text.
\par\endgroup
\ifTZPDarkEdition
\tzpdictionaryvisualband{D:/TZPID/kdp_workflow/build/tikz_figures/ID0153_dictionary_figure_dark.pdf}
\fi

\clearpage
\begingroup\tzpsupportsetup
\noindent\textbf{\fontsize{10.8pt}{12.8pt}\selectfont [ID 0153] Constructive Logic and Proof Certificate}\par
\noindent\textbf{\fontsize{10pt}{12pt}\selectfont Gravitational Stress--Energy Contribution}\par
\vspace{0.45em}
\begingroup
\footnotesize
\setlength{\parskip}{0.22em}
\setlength{\parindent}{0pt}
\noindent\textbf{Curry--Howard--Lambek Interpretation}\par
Under the Curry--Howard--Lambek correspondence, this registry entry is read as a typed proof
object: the stated source condition supplies the proposition, the derivation supplies the
morphism, and the proof certificate records the machine handles needed to check the obligation
in the formal registry.

\vspace{0.45em}
\noindent\textbf{CHL Proof}\par
\textbf{Statement:} the gravitational sector is written as the curvature-side source term that matches the
stress bookkeeping used for the vacuum and matter sectors

\textbf{Logic Validation:}\\
\textbf{Proposition:} ``Gravitational Stress--Energy Contribution is valid when the total tensorial quantity is
conserved under the stated constraint.''\\
\textbf{Proof:} The source statement identifies a conserved balance law. Reading 0153 as a proposition,
validity follows by requiring the full stress, flux, or current contribution to be
included before the divergence or balance condition is evaluated.

\textbf{VERDICT:} \ensuremath{\checkmark}\ \textbf{TRUE} --- conserved total-structure validation

\textbf{Formal Status:} \ensuremath{\checkmark}\ THEORETICAL -- source-backed CHL proof obligation

\textbf{Physical Status:} THEORETICAL -- requires model-specific empirical interpretation
\par\vspace{0.45em}
\noindent{\tiny\textbf{Isabelle/HOL.} \tzpplaincode{physics\_validity\_from\_model, external\_validity\_package, registry\_number\_matches, equation\_count\_matches}}\par
\ifTZPDarkEdition
\tzpproofvisualband{D:/TZPID/kdp_workflow/build/tikz_figures/ID0153_proof_certificate_figure_dark.pdf}
\fi
\endgroup
\endgroup
